% GENERATED FILE. Edit canonical lessons, then run npm run generate:books.
% canonical-source-hash: fnv1a64:f29f48955675a16d
% canonical-lessons: TA-C89-sadness, TA-C89-worry, TA-C89-nervousness, TA-C89-calm, TA-C89-shivering

\chapter{Sad, Worried, Calm}
\label{ch:ta-sad-worried-calm}
\hlchaptermodality{89}

\begin{chapteropening}
\textbf{By the end of this chapter:} \emph{I can say I am sad, worried, nervous or shivering, and ask someone to stay calm.}

Complete the last lesson of chapter 89: i can say I am sad, worried, nervous or shivering, and ask someone to stay calm.
\end{chapteropening}

\section[cōkam]{\ta{சோகம்} (cōkam) — sadness, grief}

\label{lesson:TA-C89-sadness}



\begin{quote}
Before the new one: say the Tamil for health, then the Tamil for sweat.
\end{quote}

\subsection*{You'll want to know: \ta{சோகம்}}
\textbf{\ta{சோகம்}} (\emph{cōkam}) — \textquotedblleft{}sadness, grief\textquotedblright{}.

\textbf{\ta{சோகமாக} \ta{இருக்கிறேன்}} — \textquotedblleft{}I am sad\textquotedblright{}. The \textbf{-\ta{ஆக}} ending turns a feeling into how you are.

\begin{grammarlens}[title={What you've built}]
The first feeling.
\end{grammarlens}

\subsection*{Your turn}
\begin{itemize}
\raggedright
  \item \emph{Say it:} \emph{cōkam}
  \item \emph{Say it:} \emph{cōkam}, once more
  \item \emph{Say it:} say \emph{cōkamāka irukkiṟēṉ}
  \item \emph{Recall:} say the Tamil for health, then the Tamil for sweat, then say \emph{cōkam} again
\end{itemize}

\begin{tcolorbox}[breakable,colback=teal!4,colframe=teal!35!black,title={Before you move on}]
What does \ta{சோகம்} mean? (\textquotedblleft{}Sadness, grief\textquotedblright{}.) Say it once more.
\end{tcolorbox}
\section[kavalai]{\ta{கவலை} (kavalai) — worry, anxiety}

\label{lesson:TA-C89-worry}



\begin{quote}
Before the new one: say the Tamil for sweat, then the Tamil for sadness.
\end{quote}

\subsection*{You'll want to know: \ta{கவலை}}
\textbf{\ta{கவலை}} (\emph{kavalai}) — \textquotedblleft{}worry\textquotedblright{}.

\textbf{\ta{கவலை} \ta{வேண்டாம்}} — \textquotedblleft{}no need to worry\textquotedblright{}, literally \textquotedblleft{}worry is not wanted\textquotedblright{}.

\begin{grammarlens}[title={What you've built}]
A worry, and a way to wave it off.
\end{grammarlens}

\subsection*{Your turn}
\begin{itemize}
\raggedright
  \item \emph{Say it:} \emph{kavalai}
  \item \emph{Say it:} \emph{kavalai}, once more
  \item \emph{Say it:} say \emph{kavalai vēṇṭām}
  \item \emph{Recall:} say the Tamil for sweat, then the Tamil for sadness, then say \emph{kavalai} again
\end{itemize}

\begin{tcolorbox}[breakable,colback=teal!4,colframe=teal!35!black,title={Before you move on}]
What does \ta{கவலை} mean? (\textquotedblleft{}Worry, anxiety\textquotedblright{}.) Say it once more.
\end{tcolorbox}
\section[pataṟṟam]{\ta{பதற்றம்} (pataṟṟam) — nervousness, tension}

\label{lesson:TA-C89-nervousness}



\begin{quote}
Before the new one: say the Tamil for sadness, then the Tamil for worry.
\end{quote}

\subsection*{You'll want to know: \ta{பதற்றம்}}
\textbf{\ta{பதற்றம்}} (\emph{pataṟṟam}) — \textquotedblleft{}nervousness, tension, a flutter\textquotedblright{}.

\textbf{\ta{பதற்றமாக} \ta{இருக்கிறது}} — \textquotedblleft{}I feel nervous\textquotedblright{}.

\begin{grammarlens}[title={What you've built}]
The opposite of calm.
\end{grammarlens}

\subsection*{Your turn}
\begin{itemize}
\raggedright
  \item \emph{Say it:} \emph{pataṟṟam}
  \item \emph{Say it:} \emph{pataṟṟam}, once more
  \item \emph{Say it:} say \emph{pataṟṟamāka irukkiṟatu}
  \item \emph{Recall:} say the Tamil for sadness, then the Tamil for worry, then say \emph{pataṟṟam} again
\end{itemize}

\begin{tcolorbox}[breakable,colback=teal!4,colframe=teal!35!black,title={Before you move on}]
What does \ta{பதற்றம்} mean? (\textquotedblleft{}Nervousness, tension\textquotedblright{}.) Say it once more.
\end{tcolorbox}
\section[amaiti]{\ta{அமைதி} (amaiti) — calm, peace, quiet}

\label{lesson:TA-C89-calm}



\begin{quote}
Before the new one: say the Tamil for worry, then the Tamil for nervousness.
\end{quote}

\subsection*{You'll want to know: \ta{அமைதி}}
\textbf{\ta{அமைதி}} (\emph{amaiti}) — \textquotedblleft{}calm, peace\textquotedblright{}, and \textquotedblleft{}quiet\textquotedblright{}.

\textbf{\ta{அமைதியாக} \ta{இருங்கள்}} — \textquotedblleft{}please stay calm\textquotedblright{}, or in a library, \textquotedblleft{}please be quiet\textquotedblright{}.

\begin{grammarlens}[title={What you've built}]
Calm, and a polite request for it.
\end{grammarlens}

\subsection*{Your turn}
\begin{itemize}
\raggedright
  \item \emph{Say it:} \emph{amaiti}
  \item \emph{Say it:} \emph{amaiti}, once more
  \item \emph{Say it:} say \emph{amaitiyāka iruṅkaḷ}
  \item \emph{Recall:} say the Tamil for worry, then the Tamil for nervousness, then say \emph{amaiti} again
\end{itemize}

\begin{tcolorbox}[breakable,colback=teal!4,colframe=teal!35!black,title={Before you move on}]
What does \ta{அமைதி} mean? (\textquotedblleft{}Calm, peace, quiet\textquotedblright{}.) Say it once more.
\end{tcolorbox}
\section[naṭukkam]{\ta{நடுக்கம்} (naṭukkam) — shivering, trembling}

\label{lesson:TA-C89-shivering}



\begin{quote}
Before the new one: say the Tamil for nervousness, then the Tamil for calm.
\end{quote}

\subsection*{You'll want to know: \ta{நடுக்கம்}}
\textbf{\ta{நடுக்கம்}} (\emph{naṭukkam}) — \textquotedblleft{}shivering, trembling\textquotedblright{}, from cold, fever or fear.

\textbf{\ta{குளிரில்} \ta{நடுக்கம்}} — \textquotedblleft{}shivering in the cold\textquotedblright{}.

\begin{grammarlens}[title={What you've built}]
Sad, worried, nervous, calm, shivering.
\end{grammarlens}

\subsection*{Your turn}
\begin{itemize}
\raggedright
  \item \emph{Say it:} \emph{naṭukkam}
  \item \emph{Say it:} \emph{naṭukkam}, once more
  \item \emph{Say it:} say \emph{naṭukkam}
  \item \emph{Recall:} say the Tamil for nervousness, then the Tamil for calm, then say \emph{naṭukkam} again
\end{itemize}

\begin{tcolorbox}[breakable,colback=teal!4,colframe=teal!35!black,title={Before you move on}]
What does \ta{நடுக்கம்} mean? (\textquotedblleft{}Shivering, trembling\textquotedblright{}.) Say it once more.
\end{tcolorbox}
